\documentclass[12pt]{amsart}
\usepackage[margin=1in]{geometry}
\pagestyle{plain}

\usepackage{amsmath,amssymb}
\usepackage{enumerate}
\usepackage[shortlabels]{enumitem}
\usepackage{graphicx}
\usepackage{xcolor}
\usepackage[T1]{fontenc}
\usepackage[parfill]{parskip}
\renewcommand{\baselinestretch}{1.1}

% Custom commands for clarity
\newcommand{\correct}[1]{{\color{blue}\textbf{CORRECT: #1}}}
\newcommand{\partcred}[1]{{\color{red}\textit{Partial Credit: #1}}}
\newcommand{\hint}{\textbf{Key Insight: }}
\newcommand{\trap}{\textbf{Common Trap: }}

\title{ESE 2030 QUIZZAM 1 : Solution Guide\\Fall 2025}
\author{Solution Key with Explanations}

\begin{document}
\maketitle

\section*{Overview}
This solution guide provides detailed explanations for each problem on Quiz 1. Pay special attention to:
\begin{itemize}
\item Why the correct answer is right
\item Common misconceptions that lead to wrong answers
\item Partial credit opportunities where applicable
\item Key concepts being tested
\end{itemize}

\hrulefill

\section*{Problem 1: Dimension of Polynomial Spaces}
\textbf{Topic:} Polynomial spaces, dimension, bases

\textbf{Problem:} The dimension of $\mathcal{P}_3$ (polynomials of degree at most 3).

\correct{(C) 4}

\textbf{Explanation:} The space $\mathcal{P}_3$ consists of all polynomials of the form:
$$p(x) = a_0 + a_1x + a_2x^2 + a_3x^3$$

A standard basis is $\{1, x, x^2, x^3\}$, which has 4 elements. Therefore, $\dim(\mathcal{P}_3) = 4$.

\hint The subscript in $\mathcal{P}_n$ denotes the maximum degree, not the dimension. The dimension is always $n+1$.

\trap Students often confuse "degree at most 3" with "exactly degree 3" or think the dimension equals the maximum degree.

\hrulefill

\section*{Problem 2: Triangular Systems}
\textbf{Topic:} Triangular systems, forward substitution, computational efficiency

\correct{(C) One should solve by forward substitution}

\partcred{(A) Row reduction is possible but inefficient}

\textbf{Explanation:} For a lower triangular system $L\mathbf{x} = \mathbf{b}$:
\begin{itemize}
\item Since $L$ has nonzero diagonal entries, it's invertible, so there's always a unique solution
\item Forward substitution is the most efficient method:
  \begin{enumerate}
  \item Solve $2x_1 = b_1$ to get $x_1 = b_1/2$
  \item Substitute into second equation: $x_1 + 3x_2 = b_2$
  \item Continue forward through the system
  \end{enumerate}
\item Backward substitution is for upper triangular systems
\end{itemize}

\hint Forward = Lower, Backward = Upper (alphabetical order helps remember!)

\hrulefill

\section*{Problem 3: Kernel Dimension}
\textbf{Topic:} Linear transformations, kernel, rank-nullity theorem

\correct{(C) 2}

\partcred{(B) 1 - shows understanding but miscalculation}

\textbf{Explanation:} 
Observe that all three rows of $A$ are scalar multiples:
\begin{align}
\text{Row 2} &= 2 \times \text{Row 1}\\
\text{Row 3} &= -1 \times \text{Row 1}
\end{align}

Therefore:
\begin{itemize}
\item $\text{rank}(A) = 1$ (only one independent row)
\item By the Rank-Nullity Theorem: $\dim(V) = \text{rank}(T) + \dim(\ker(T))$
\item So: $3 = 1 + \dim(\ker(T))$
\item Therefore: $\dim(\ker(T)) = 2$
\end{itemize}

\hint When all rows/columns are proportional, the rank is 1.

\hrulefill

\section*{Problem 4: Vector Space Examples}
\textbf{Topic:} Vector spaces, axioms, closure properties

\correct{(B) Continuous functions with $f(0) = f(1) = 0$}

\textbf{Explanation:} Let's check each option:

\textbf{(A) Matrices with det = 1:} Fails closure under addition. If $\det(A) = \det(B) = 1$, generally $\det(A+B) \neq 1$.

\textbf{(B) Functions with $f(0) = f(1) = 0$:} ✓ This works!
\begin{itemize}
\item Zero function satisfies the condition
\item If $f(0) = g(0) = 0$ and $f(1) = g(1) = 0$, then $(f+g)(0) = 0$ and $(f+g)(1) = 0$
\item If $f(0) = f(1) = 0$, then $(cf)(0) = cf(0) = 0$ and $(cf)(1) = cf(1) = 0$
\end{itemize}

\textbf{(C) Vectors with $x+y+z=1$:} Not a subspace; doesn't contain zero vector.

\textbf{(D) Idempotent matrices:} Fails closure. If $A^2 = A$ and $B^2 = B$, generally $(A+B)^2 \neq A+B$.

\textbf{(E) First quadrant:} Fails closure under negative scalar multiplication.

\hint Always check: (1) Contains zero? (2) Closed under addition? (3) Closed under scalar multiplication?

\hrulefill

\section*{Problem 5: Differentiation Operator}
\textbf{Topic:} Linear operators, rank and nullity

\correct{(A) rank = 3, nullity = 1}

\textbf{Explanation:} The differentiation operator $\mathcal{D}: \mathcal{P}_3 \to \mathcal{P}_2$ works as:
$$\mathcal{D}(a_0 + a_1x + a_2x^2 + a_3x^3) = a_1 + 2a_2x + 3a_3x^2$$

\begin{itemize}
\item The image is all of $\mathcal{P}_2$ (any quadratic can be obtained), so $\text{rank} = \dim(\mathcal{P}_2) = 3$
\item The kernel consists of constant polynomials (derivatives equal zero), so $\dim(\ker(\mathcal{D})) = 1$
\item Verify: $\dim(\mathcal{P}_3) = 4 = 3 + 1$ ✓
\end{itemize}

\hint The derivative of constants is zero, so the kernel is the space of constant polynomials.

\hrulefill

\section*{Problem 6: Elementary Matrices}
\textbf{Topic:} Elementary matrices, determinants, invertibility

\correct{(C) $\det(E) \neq 0$}

\partcred{(E) Only true for row swaps}

\textbf{Explanation:} Elementary matrices correspond to three types of row operations:
\begin{enumerate}
\item Row swaps: $\det(E) = -1$
\item Row scaling by $c$: $\det(E) = c \neq 0$
\item Adding multiple of one row to another: $\det(E) = 1$
\end{enumerate}

All elementary matrices are invertible (operations can be undone), so $\det(E) \neq 0$ always.

\trap Not all elementary matrices have $\det = \pm 1$; row scaling can give any nonzero determinant.

\hrulefill

\section*{Problem 7: Coimage and Cokernel}
\textbf{Topic:} Quotient spaces, fundamental theorem

\correct{(C) $\text{coim } T = V/\ker T$ and $\text{coker } T = W/\text{im } T$}

\textbf{Explanation:} By definition:
\begin{itemize}
\item \textbf{Coimage:} The quotient of the domain by the kernel: $V/\ker T$
  \begin{itemize}
  \item This "mods out" the part that maps to zero
  \item Results in equivalence classes of vectors that map to the same output
  \end{itemize}
\item \textbf{Cokernel:} The quotient of the codomain by the image: $W/\text{im } T$
  \begin{itemize}
  \item This captures "what's missing" from the image
  \item Measures the failure of surjectivity
  \end{itemize}
\end{itemize}

\hint Remember: "co-" operations involve quotients of the corresponding spaces.

\hrulefill

\section*{Problem 8: Permutation Matrices}
\textbf{Topic:} Permutation matrices, cyclic permutations

\correct{(B) $P^3 = I$}

\partcred{(D) $P^{-1} = P$ - true for 2-cycles but not this 3-cycle}

\textbf{Explanation:} This permutation matrix cycles rows: $1 \to 2 \to 3 \to 1$.

Computing powers:
\begin{itemize}
\item $P$ maps: row 1 → row 2, row 2 → row 3, row 3 → row 1
\item $P^2$ maps: row 1 → row 3, row 2 → row 1, row 3 → row 2
\item $P^3$ maps: row 1 → row 1, row 2 → row 2, row 3 → row 3 (identity!)
\end{itemize}

For this 3-cycle: $P^{-1} = P^2 \neq P$ (unlike 2-cycles where $P^{-1} = P$).

\hint For an $n$-cycle permutation, $P^n = I$.

\hrulefill

\section*{Problem 9: Linear Independence in $\mathbb{R}^3$}
\textbf{Topic:} Linear independence, span

\correct{(C) Linearly dependent and does not span $\mathbb{R}^3$}

\textbf{Explanation:} Notice that:
$$\begin{pmatrix}1\\3\\1\end{pmatrix} = \begin{pmatrix}1\\2\\0\end{pmatrix} + \begin{pmatrix}0\\1\\1\end{pmatrix}$$

This shows linear dependence. Since we have only 2 independent vectors in $\mathbb{R}^3$, the set cannot span the entire space.

\trap Three vectors in $\mathbb{R}^3$ aren't automatically a basis—check for independence!

\hrulefill

\section*{Problem 10: Subspace Operations}
\textbf{Topic:} Subspaces, intersection and union

\correct{(B) $W_1 \cap W_2$}

\textbf{Explanation:} 
\begin{itemize}
\item \textbf{(A) Union:} Generally NOT a subspace. Example: $x$-axis ∪ $y$-axis in $\mathbb{R}^2$
\item \textbf{(B) Intersection:} Always a subspace! ✓
  \begin{itemize}
  \item Contains zero (in both subspaces)
  \item Closed under addition and scalar multiplication
  \end{itemize}
\item \textbf{(C) Set difference:} Not closed under scalar multiplication
\item \textbf{(D) Complement of sum:} Not closed under operations
\item \textbf{(E) Quotient spaces:} These ARE subspaces but of $V/W_i$, not of $V$
\end{itemize}

\hint Intersection preserves all subspace properties; union typically doesn't.

\hrulefill

\section*{Problem 11: Symmetric Matrices}
\textbf{Topic:} Dimension of matrix subspaces

\correct{(D) 6}

\partcred{(B) 3 - counts only diagonal}

\textbf{Explanation:} A $3 \times 3$ symmetric matrix has the form:
$$\begin{bmatrix} a & b & c \\ b & d & e \\ c & e & f \end{bmatrix}$$

Count the independent parameters:
\begin{itemize}
\item Diagonal: 3 parameters $(a, d, f)$
\item Upper triangle: 3 parameters $(b, c, e)$
\item Total: 6 dimensions
\end{itemize}

Formula: For $n \times n$ symmetric matrices, $\dim = \frac{n(n+1)}{2}$

\hrulefill

\section*{Problem 12: Dimension Constraints}
\textbf{Topic:} Rank-nullity theorem, dimension bounds

\correct{(C) The dimension of $\ker(T)$ is at least 2}

\textbf{Explanation:} By the Rank-Nullity Theorem:
$$\dim(V) = \dim(\ker T) + \dim(\text{im } T)$$

Since $\text{im}(T) \subseteq W$, we have $\dim(\text{im } T) \leq \dim(W) = 3$.

Therefore:
$$5 = \dim(\ker T) + \dim(\text{im } T) \leq \dim(\ker T) + 3$$

So $\dim(\ker T) \geq 2$.

\hint When domain dimension exceeds codomain dimension, there must be a nontrivial kernel.

\hrulefill

\section*{Problem 13: LU Decomposition and Nonsingularity}
\textbf{Topic:} LU decomposition, triangular matrices, determinants

\correct{(E) $A$ is nonsingular if and only if all diagonal entries of $U$ are nonzero}

\partcred{(B) Close but says "last" instead of "all"; (D) Wrong—signs don't matter}

\textbf{Explanation:} 
\begin{itemize}
\item In standard LU decomposition, $L$ has 1's on the diagonal (by convention)
\item Therefore, $L$ is always invertible
\item Since $A = LU$ and $\det(A) = \det(L)\det(U) = \det(U)$
\item For triangular matrices: nonsingular ⟺ all diagonal entries nonzero
\item Thus: $A$ nonsingular ⟺ $U$ nonsingular ⟺ all diagonal entries of $U$ are nonzero
\end{itemize}

\hrulefill

\section*{Problem 14: LU Properties}
\textbf{Topic:} LU decomposition fundamentals

\correct{(A) The matrix $L$ is invertible}

\textbf{Explanation:} By convention in LU decomposition:
\begin{itemize}
\item $L$ is lower triangular with 1's on the diagonal
\item This guarantees $\det(L) = 1 \neq 0$
\item Therefore, $L$ is always invertible
\item $U$ may have zeros on diagonal (singular case)
\item The other options depend on whether $A$ is singular
\end{itemize}

\hint The "L" in LU always has unit diagonal, making it invertible by construction.

\hrulefill

\section*{Problem 15: Basis Extension}
\textbf{Topic:} Basis, unique representation

\correct{(B) $\mathbf{v}_5$ can be uniquely written as a linear combination of $\mathbf{v_1, v_2, v_3, v_4}$}

\textbf{Explanation:} Since $\{\mathbf{v_1, v_2, v_3, v_4}\}$ is a basis for the 4-dimensional space $V$:
\begin{itemize}
\item Every vector in $V$ (including $\mathbf{v}_5$) can be written as a linear combination of the basis
\item This representation is unique (fundamental property of bases)
\item $\mathbf{v}_5$ need not be zero (option A is wrong)
\item The set $S$ is linearly dependent since it has 5 vectors in a 4D space
\end{itemize}

\hint In an $n$-dimensional space, any set of more than $n$ vectors must be linearly dependent.

\hrulefill

\section*{Problem 16: Pivots and Solutions}
\textbf{Topic:} Row reduction, pivot analysis

\correct{(B) The system has infinitely many solutions if it has any solution at all}

\partcred{(E) Ambiguity about pivot location}

\textbf{Explanation:} With 2 pivots in a $3 \times 3$ system:
\begin{itemize}
\item One row lacks a pivot (becomes a zero row or inconsistent row)
\item If the third row is $[0\ 0\ 0\ |\ b_3]$ with $b_3 \neq 0$: inconsistent
\item If the third row is $[0\ 0\ 0\ |\ 0]$: one free variable → infinitely many solutions
\item The system either has no solutions or infinitely many (never unique)
\end{itemize}

\hint Number of free variables = number of columns minus number of pivots in coefficient matrix.

\hrulefill

\section*{Problem 17: Non-Linear Transformation}
\textbf{Topic:} Linear transformation properties

\correct{(D) $F(x,y) = (x+1, y, x-y)$}

\textbf{Explanation:} For a transformation to be linear, it must satisfy $T(\mathbf{0}) = \mathbf{0}$.

Check $F$: $F(0,0) = (0+1, 0, 0-0) = (1, 0, 0) \neq (0, 0, 0)$

This immediately disqualifies $F$. Additionally, $F$ fails additivity:
$$F(x_1+x_2, y_1+y_2) = (x_1+x_2+1, y_1+y_2, x_1+x_2-y_1-y_2)$$
$$\neq F(x_1,y_1) + F(x_2,y_2) = (x_1+1, y_1, x_1-y_1) + (x_2+1, y_2, x_2-y_2)$$

All other options preserve the zero vector and satisfy linearity.

\trap The "+1" is the killer—any additive constant breaks linearity.

\hrulefill

\section*{Problem 18: Fundamental Theorem of Linear Algebra}
\textbf{Topic:} Isomorphism between coimage and image

\correct{(E) None of the above}

\textbf{Explanation:} The Fundamental Theorem of Linear Algebra states:
$$\text{coim}(T) \cong \text{im}(T)$$

That is, the coimage $V/\ker(T)$ is isomorphic to the image of $T$.

None of the given options correctly state this relationship. This tests precise understanding of the theorem's statement.

\hint The key isomorphism is between what's "left after modding out the kernel" and what actually gets mapped to.

\hrulefill

\section*{Problem 19: Matrix Transformations}
\textbf{Topic:} Correspondence between matrices and linear maps

\correct{(C) The null space of $A$ is the kernel of $T_A$}

\partcred{(B) Column space equals image, not coimage}

\textbf{Explanation:} By definition:
\begin{itemize}
\item The null space of $A$ = $\{\mathbf{x} : A\mathbf{x} = \mathbf{0}\}$ = $\ker(T_A)$ ✓
\item The column space of $A$ = image of $T_A$ (not coimage)
\item "Nonsingular" only makes sense for square matrices
\item Options (A), (D), (E) incorrectly relate injectivity/surjectivity to singularity
\end{itemize}

\hrulefill

\section*{Problem 20: Spanning Definition}
\textbf{Topic:} Definition of span

\correct{(E) $S$ spans $V$}

\partcred{(C) True but not "by definition"}

\textbf{Explanation:} The problem statement exactly describes the definition of "$S$ spans $V$."

\begin{itemize}
\item This doesn't imply $S$ is a basis (need linear independence too)
\item This doesn't imply $S$ is linearly independent
\item While (C) is true (can't span with fewer than $\dim(V)$ vectors), it's a theorem, not the definition
\end{itemize}

\hint "Spans" means exactly "every vector can be written as a linear combination."

\hrulefill

\section*{Problem 21: Polynomial Independence}
\textbf{Topic:} Linear independence in $\mathcal{P}_2$

\correct{(A) Linearly independent and spans $\mathcal{P}_2$}

\textbf{Explanation:} Check linear independence by setting up:
$$c_1(1+x) + c_2(x+x^2) + c_3(1+x^2) = 0$$

This gives the system:
\begin{align}
c_1 + c_3 &= 0 \quad \text{(constant term)}\\
c_1 + c_2 &= 0 \quad \text{(coefficient of } x\text{)}\\
c_2 + c_3 &= 0 \quad \text{(coefficient of } x^2\text{)}
\end{align}

The only solution is $c_1 = c_2 = c_3 = 0$, so the set is linearly independent.

Since we have 3 independent vectors in the 3-dimensional space $\mathcal{P}_2$, they form a basis.

\hrulefill

\section*{Problem 22: Cokernel Dimension}
\textbf{Topic:} Cokernel, quotient spaces

\correct{(D) $\mathbb{R}^3$}

\textbf{Explanation:} 
Given: $T: \mathbb{R}^5 \to \mathbb{R}^6$ with $\dim(\ker(T)) = 2$

By Rank-Nullity: $\dim(\text{im}(T)) = 5 - 2 = 3$

The cokernel is $\text{coker}(T) = \mathbb{R}^6/\text{im}(T)$

Therefore: $\dim(\text{coker}(T)) = 6 - 3 = 3$

So coker$(T) \cong \mathbb{R}^3$.

\hint The cokernel measures "what's missing" from the image to fill the codomain.

\hrulefill

\section*{Problem 23: Quotient Space Elements}
\textbf{Topic:} Equivalence classes, quotient by subspaces

\correct{(C) $(-1, -2, 2)^T$}

\textbf{Explanation:} Two vectors are equivalent if their difference lies in $W$ (the $xy$-plane).

For $\mathbf{v} = (3, 5, 2)^T$, check each option:

\begin{itemize}
\item (A): $(3,5,2) - (3,5,0) = (0,0,2) \notin W$ (has nonzero $z$-component)
\item (B): $(3,5,2) - (6,10,4) = (-3,-5,-2) \notin W$
\item (C): $(3,5,2) - (-1,-2,2) = (4,7,0) \in W$ ✓ (zero $z$-component)
\item (D): $(3,5,2) - (0,0,0) = (3,5,2) \notin W$
\end{itemize}

\hint Vectors are equivalent mod $W$ if they differ only by an element of $W$.

\hrulefill

\section*{Problem 24: Elementary Matrix Order}
\textbf{Topic:} Composition of row operations

\correct{(A) $B = E_2 E_1 A$}

\textbf{Explanation:} Elementary matrices apply from right to left:
\begin{enumerate}
\item First operation: $E_1 A$ swaps rows 1 and 2
\item Second operation: $E_2(E_1 A)$ adds 2 times (new) row 1 to row 3
\end{enumerate}

The key insight: After swapping, "row 1" in the second operation refers to what was originally row 2.

\trap Matrix multiplication is right-to-left, but we describe operations left-to-right!

\hrulefill

\section*{Problem 25: Subspace Violations}
\textbf{Topic:} Subspace axioms

\correct{(E) All three: (A), (B), and (C)}

\partcred{(D) Identifies two violations correctly}

\textbf{Explanation:} The set $S = \{(x,y,z) : x+y+z = 1\}$ is an affine space (shifted plane).

Check all three properties:
\begin{enumerate}
\item \textbf{Zero vector:} $(0,0,0) \notin S$ since $0+0+0 = 0 \neq 1$ ✗
\item \textbf{Addition:} If $\mathbf{u}, \mathbf{v} \in S$, then sum of coordinates of $\mathbf{u}+\mathbf{v}$ is $1+1 = 2 \neq 1$ ✗
\item \textbf{Scalar multiplication:} If $\mathbf{u} \in S$ and $c \neq 1$, then sum of coordinates of $c\mathbf{u}$ is $c \cdot 1 = c \neq 1$ ✗
\end{enumerate}

All three properties fail!

\hint Any affine space (not through origin) fails all subspace properties.

\hrulefill

\section*{Study Tips}

\begin{enumerate}
\item \textbf{Conceptual Understanding:} Focus on definitions and theorems, not just computation
\item \textbf{Check Special Cases:} Always verify with simple examples (zero vector, identity matrix, etc.)
\item \textbf{Watch for Keywords:} "Must be," "always," "if and only if" signal precise mathematical statements
\item \textbf{Dimension Counting:} Many problems reduce to careful dimension arithmetic
\item \textbf{Subspace Tests:} Always check all three properties systematically
\end{enumerate}

\end{document}